% Згенеровано run_all.m. Абзаци з результатами перевірки в GNU Octave через GLPK.
% Початковий текст роботи не змінено; ці абзаци подано окремо.

\paragraph{Середовище та параметри розв’язувача.}
Обчислення виконано в GNU Octave 8.4.0 функцією \texttt{glpk}, яка викликає бібліотеку GLPK 5.0 (Ubuntu 24.04.5 LTS). Усі задачі розв’язано двофазним прямим симплекс-методом без пресолвера (\texttt{lpsolver}=1, \texttt{dual}=1, \texttt{presol}=0). Використано рівноважне масштабування (\texttt{scale}=16), вибір змінної за найкрутішим ребром (\texttt{price}=34) і тест відношень Гарріса (\texttt{rtest}=34). Допуски розв’язувача: \texttt{tolbnd}$=10^{-9}$, \texttt{toldj}$=10^{-9}$, \texttt{tolpiv}$=10^{-10}$. Запуск вважали успішним лише за коду завершення~0 і статусу~5 (оптимальний розв’язок). Цю умову виконано в усіх 294 запусках. Носій визначали умовою $x_i^*>10^{-9}$, активні рядки --- умовою $|s_j|\le10^{-9}\max\{1,|b_j|\}$.

\paragraph{Вихідні дані.}
Матрицю 1939 року взято з файлу \texttt{stigler.mod} архіву GLPK~4.65 без змін: контрольна сума SHA-256 файлу (\texttt{b0930f90e9b9ddee\ldots}) збігається із сумою оригіналу. Файл містить 77 продуктів і дев’ять показників у порядку: енергія, білок, кальцій, залізо, вітамін~A, тіамін, рибофлавін, ніацин, вітамін~C. Коефіцієнти задано на один долар витрат, ціни --- в центах за торгову одиницю. Для 58 продуктів із ваговою торговою одиницею перевірено стовпець їстівної маси на один долар: у 57 випадках він узгоджується з ціною та масою торгової одиниці. Для продукту \texttt{cornsirup} файл містить 4996~г, тоді як торгова одиниця 24 унції за 13{,}7~цента дає 4966{,}4~г; цей стовпець до моделі не входить. 

Матрицю 1944 року утворюють 14 рядків таблиці~B на с.~307 статті Стіглера. Таблиця~B не містить стовпця торгових одиниць: для продуктів №~1--74 одиниці взято з таблиці~A, для №~78--80 --- з приміток до таблиці~B. Порожні клітинки передано нулями. Звірка файлу \texttt{stigler1944.csv} зі сканом сторінки підтвердила 139 зі 140 клітинок (ціна та дев’ять показників у кожному рядку). Для 11 продуктів, які є в таблиці~A, перевірено співвідношення $a^{1944}_{ji}=a^{1939}_{ji}p^{1939}_i/p^{1944}_i$. Відхилення не перевищує 0{,}53 одиниці останнього друкованого знака в усіх клітинках, крім однієї. Виняток --- білок продукту №~3: у файлі 398~г, за таблицею~A --- 393{,}25~г; остання цифра на скані читається неоднозначно (3 або 8). Заміна 398 на 393 не змінює $C^*$, $x^*$, $u^*$, зведені вартості та значення $m_j$: цей продукт не входить до носія, а множник білкового рядка дорівнює нулю. Для продуктів №~78 і №~80 їстівна маса на один долар із приміток (4647 і 2071~г) узгоджується з ціною та торговою одиницею (4647{,}5 і 2071{,}2~г).

\paragraph{Пряма і двоїста задачі.}
Для кожної матриці окремо розв’язано пряму задачу $\min\{\mathbf1^\top x:\ Ax\ge b,\ x\ge0\}$ і двоїсту задачу $\max\{b^\top u:\ A^\top u\le\mathbf1,\ u\ge0\}$. Функція \texttt{glpk} повертає множники рядків як похідні оптимального значення за правими частинами. Тому для прямої задачі з рядками «$\ge$» множники збігаються з $u^*\ge0$, а для двоїстої задачі з рядками «$\le$» --- з $x^*\ge0$; зміна знака не потрібна. Множники прямої задачі відрізняються від розв’язку двоїстої не більше ніж на $7{,}63\cdot10^{-17}$ (1939) і $1{,}73\cdot10^{-17}$ (1944).

Одержано $C^*_{1939}=0{,}108662278206757$ і $C^*_{1944}=0{,}164110826270237$ дол./день. За 365 днів це 39{,}661732 і 59{,}900452 дол., за $365{,}25$ дня --- 39{,}688897 і 59{,}941479 дол. У табл.~\ref{tab:glpk-solution-1939} і \ref{tab:glpk-solution-1944} наведено оптимальні вектори $x^*$, у табл.~\ref{tab:glpk-dual-1939} і \ref{tab:glpk-dual-1944} --- надходження $Ax^*$, запаси $s=Ax^*-b$ і множники $u^*$. Повні вектори з 77 і 14 компонент та зведені вартості $r=\mathbf1-A^\top u^*$ збережено у файлах CSV. Порушення умов $Ax^*\ge b$, $x^*\ge0$, $A^\top u^*\le\mathbf1$, $u^*\ge0$ не перевищує $2{,}84\cdot10^{-14}$, двоїстий розрив --- $1{,}11\cdot10^{-16}$, добутки $|u_j^*s_j|$ і $|x_i^*r_i|$ --- $1{,}08\cdot10^{-16}$ (табл.~\ref{tab:glpk-checks}). В обох задачах додатними є п’ять компонент $x^*$ і п’ять компонент $u^*$. Активні рядки: енергія, кальцій, вітамін~A, вітамін~B$_2$, вітамін~C.

\paragraph{Фізичні кількості.}
Кількість торгових одиниць обчислено за формулою $z_i=100x_i^*/p_i$, де $p_i$ --- ціна торгової одиниці в центах. У табл.~\ref{tab:glpk-units-1939} і \ref{tab:glpk-units-1944} наведено добові кількості та річні кількості за 365 і $365{,}25$ дня. Торгова одиниця пшеничного борошна дорівнює 10 фунтам, борошна для млинців --- 20 унціям, тобто $1{,}25$ фунта. За $365{,}25$ дня розв’язок 1944 року містить 534{,}151 фунта пшеничного борошна і 134{,}414 фунта борошна для млинців, за 365 днів --- 533{,}786 і 134{,}322 фунта.

\paragraph{Історичний підсумок 1944 року.}
Сума вартостей у стовпці 1944 року таблиці~2 дорівнює 59{,}88 дол. Мінімум для опублікованої матриці становить 59{,}9005 дол. за 365 днів і 59{,}9415 дол. за $365{,}25$ дня. Стаття Стіглера не називає кількість днів у році; сума $24{,}50$ дол. на с.~311 відповідає 365 дням ($365\cdot3/44{,}7=24{,}50$). Тому історичний підсумок зіставлено зі значенням за 365 днів: різниця дорівнює 0{,}0205 дол., або 0{,}034\%. Історичний раціон містить ті самі п’ять продуктів, що й носій $x^*$ (табл.~\ref{tab:glpk-historic-1944}). Якщо його річні вартості поділити на 365, то за коефіцієнтами таблиці~B добове надходження становить 99{,}83\% вимоги для кальцію, 99{,}91\% вимоги для вітаміну~A, 99{,}82\% вимоги для вітаміну~B$_2$ (табл.~\ref{tab:glpk-historic-1944-nutrients}). Отже, історичний раціон не є допустимим для надрукованої матриці, і його вартість не суперечить мінімуму 59{,}9005 дол. За властивістю $V(\theta b)=\theta V(b)$ бюджет 59{,}88 дол. покриває щонайбільше 99{,}966\% усіх вимог одночасно. Коефіцієнти таблиці~B надруковано з кроком $0{,}1$ або 1. Якщо всі ненульові коефіцієнти збільшити на половину кроку, мінімальна вартість за 365 днів дорівнює 59{,}6902 дол.; якщо зменшити --- 60{,}1122 дол. Мінімум для кожної матриці, сумісної з округленням, лежить між цими значеннями, оскільки поелементне збільшення $A$ не збільшує $V$. Історичний підсумок 59{,}88 дол. належить цьому проміжку. Перевірка не відновлює неокруглені дані Стіглера.

\paragraph{Єдиність оптимального вектора.}
На множині $\{x:\ Ax\ge b,\ x\ge0,\ \mathbf1^\top x=C^*\}$ із повним обчисленим значенням $C^*$ знайдено мінімум і максимум кожної координати: 154 задачі для 1939 року і 28 задач для 1944 року. Усі задачі завершено з кодом~0 і статусом~5. Найбільша різниця між верхньою та нижньою межами координати дорівнює $5{,}95\cdot10^{-15}$ і $3{,}75\cdot10^{-16}$ дол./день (табл.~\ref{tab:glpk-uniqueness}). Це числовий результат, одержаний із допусками \texttt{tolbnd}$=$\texttt{toldj}$=10^{-9}$; він не є доведенням єдиності.

Строгий висновок одержано точною перевіркою сертифіката. Усі коефіцієнти мають щонайбільше один десятковий знак, тому $10A$ і $10b$ --- цілочислові. Для $S=\{i:x_i^*>0\}$ і $J_+=\{j:u_j^*>0\}$ у цілочисловій арифметиці без округлень обчислено визначник і приєднану матрицю для $10A_{J_+,S}$. Одержано $\det(10A^{1939}_{J_+,S})=-10457121500604640$ і $\det(10A^{1944}_{J_+,S})=5346316706566060$, тобто $\det A^{1939}_{J_+,S}=-65357009378779/625$ і $\det A^{1944}_{J_+,S}=267315835328303/5000$. Точно перевірено такі умови: $x_S^*>0$; $u_{J_+}^*>0$; $Ax^*\ge b$ з рівністю в рядках $J_+$; $A^\top u^*\le\mathbf1$ з рівністю в стовпцях $S$; $r_i>0$ для всіх $i\notin S$; $\mathbf1^\top x^*=b^\top u^*$. Найменша зведена вартість поза носієм дорівнює 0{,}04366642 для 1939 року і 0{,}09656077 для 1944 року. Отже, для обох табличних матриць оптимальний вектор $x^*$ єдиний. Запаси в рядках поза $J_+$ строго додатні, тому двоїстий оптимум $u^*$ також єдиний. Висновок стосується зазначених десяткових коефіцієнтів.

\paragraph{Чутливість до правих частин.}
Кожну з дев’яти вимог збільшено окремо на $1\%$ за незмінних інших. В усіх 18 випадках фактичний приріст вартості збігається з $0{,}01\,b_ju_j^*$; найбільша абсолютна різниця дорівнює $2{,}78\cdot10^{-17}$ дол. (табл.~\ref{tab:glpk-sensitivity}). Точна перевірка показала, що в кожному з цих випадків базис $(S,J_+)$ залишається оптимальним. Тому рівність $\Delta C^*=0{,}01\,b_ju_j^*$ для цих 18 збільшень є точною. Результат не поширюється на довільні зміни правих частин. За одночасного збільшення всіх вимог на $1\%$ одержано 0{,}109748900989 і 0{,}165751934533 дол./день. Приріст дорівнює $0{,}01\,C^*$ з різницею не більше ніж $1{,}11\cdot10^{-17}$ дол. і збігається із сумою дев’яти окремих приростів.

\paragraph{Неактивність і надлишковість.}
Для кожного неактивного рядка обчислено $m_j=\min\{A_{j\cdot}x:\ A_{k\cdot}x\ge b_k\ (k\ne j),\ x\ge0\}$ (табл.~\ref{tab:glpk-redundancy}). В усіх 8 випадках $m_j<b_j$, тому жодне неактивне обмеження не є надлишковим. Вилучення будь-якого одного рядка з $u_j^*=0$ і всіх чотирьох таких рядків одночасно змінює мінімальну вартість не більше ніж на $2{,}08\cdot10^{-16}$ дол./день, а оптимальний вектор --- не більше ніж на $7{,}89\cdot10^{-17}$ дол./день.

\paragraph{Семиденна задача.}
Сума семи окремих добових мінімумів дорівнює 0{,}796714750932 дол., мінімум спільної семиденної задачі --- 0{,}796714750932 дол.; різниця дорівнює $0$ (табл.~\ref{tab:glpk-multiday-days}). Добовий мінімум дорівнює 0{,}108662278207 дол. для вимог $(3;\ 70)$, 0{,}113403160673 дол. для $(3{,}3;\ 73{,}5)$ і 0{,}119590356250 дол. для $(3{,}6;\ 77)$. У дні 1 і 4 носій складається з п’яти продуктів. У дні 2, 3, 5, 6, 7 він складається з чотирьох продуктів (без яловичої печінки), а рядок вітаміну~B$_2$ неактивний. Кожний із семи добових оптимумів єдиний (точна перевірка сертифіката). Тому сукупність добових оптимумів єдина, а середні надходження в ній визначено однозначно.

Середнє надходження заліза за сім днів у цій сукупності дорівнює 62{,}039142467~мг. Порівняння в цілих числах показує, що воно строго більше за 12 і 62{,}03914~мг та строго менше за 65{,}14113~мг. За твердженням~1 межі 12 і 62{,}03914~мг зберігають $C^*$, а межа 65{,}14113~мг збільшує її. GLPK дає $C^*=0{,}796714751$ дол. для перших двох меж і $C^*=0{,}809187146$ дол. для третьої; приріст дорівнює 0{,}012472395 дол. (табл.~\ref{tab:glpk-multiday-iron}).

Для білка на множині $D_j=\{2,3,5,6,7\}$ середня добова нижня межа дорівнює 74{,}9~г, а середнє надходження --- 158{,}661312557~г. Воно строго більше за 74{,}9 і 158{,}66131~г та строго менше за 166{,}594~г. GLPK дає $C^*=0{,}796714751$ дол. для перших двох меж і $C^*=0{,}815475724$ дол. для третьої; приріст дорівнює 0{,}018760973 дол. (табл.~\ref{tab:glpk-multiday-protein}).

Допоміжні максимізації середнього надходження виконано за бюджету $0{,}79671475+10^{-7}$ дол. Найбільше середнє надходження заліза дорівнює 62{,}039167~мг, білка на множині $D_j$ --- 158{,}661354~г. Це результати з допуском до вартості: бюджет перевищує $C^*$ на $9{,}91\cdot10^{-8}$ дол. На точній оптимальній грані $\mathbf1^\top x=C^*$ мінімум і максимум середнього надходження збігаються: 62{,}039142~мг для заліза і 158{,}661313~г для білка.

\paragraph{Порівняння з текстом роботи.}
У табл.~\ref{tab:glpk-comparison} значення з тексту зіставлено з результатами GLPK; таблиця містить 164 рядки. Для 162 значень, які не залежать від розв’язувача, результат GLPK після округлення до наведеної в тексті кількості знаків збігається з текстом у 162 випадках. Числових розбіжностей не виявлено. Значення $1{,}23\cdot10^{-15}$ і $2{,}02\cdot10^{-16}$ у підрозд.~2.3 є числовими похибками іншого розв’язувача; GLPK дає $5{,}95\cdot10^{-15}$ і $3{,}75\cdot10^{-16}$.

\paragraph{Запропоновані виправлення.}
\begin{enumerate}
\item Підрозд.~1.1, «Обчислювальна перевірка». Речення про \texttt{glpsol} версії~4.65, SciPy і SymPy описує попередній запуск. Для результатів цього додатка середовище таке: GNU Octave 8.4.0, бібліотека GLPK 5.0; версія~4.65 стосується лише файлу даних \texttt{stigler.mod}.
\item Підрозд.~2.3. Якщо різниці меж координат наводяться для GLPK, значення $1{,}23\cdot10^{-15}$ і $2{,}02\cdot10^{-16}$ слід замінити на $5{,}95\cdot10^{-15}$ і $3{,}75\cdot10^{-16}$.
\item Файл \texttt{stigler1944.csv}, продукт №~3, білок. Значення 398 не підтверджено однозначно; за таблицею~A очікується 393. На числа розділу ця клітинка не впливає.
\item Підрозд.~1.1, зіставлення з підсумком $59{,}88$. Текст порівнює його зі значенням за 365 днів (59{,}9005 дол.); це узгоджується зі статтею. Доцільно додати обґрунтування: стаття не називає кількість днів, але сума $24{,}50$ дол. на с.~311 відповідає 365 дням.
\end{enumerate}
